\section{Threat Model, Assumptions, and Goals}
\label{sec:threat}

\subsection{System and adversary}
The system is the master, the switch, and the SEL-751 outstation of Section~\ref{sec:bg}, exchanging
the native read and control transactions defined there. The adversary is a \emph{passive} observer on
the master-facing link: it captures traffic from a mirrored port or a tap and does not modify, inject,
or block a packet. Its goal is reconnaissance, to fingerprint the outstation and its transactions from
what the traffic exposes.

\subsection{What the adversary can and cannot read}
The adversary reads metadata: TCP and IP headers, TCP segment sizes, and arrival times. The traffic is
unencrypted, so the DNP3 payload is in fact parseable; our threat, however, uses only the two
metadata features of Section~\ref{sec:bg} (CLRT and segment shape) and does not depend on application
semantics. Following Formby et~al.~\cite{formby2016}, we distinguish three fingerprinting aims: device
identity, device type, and transaction-class inference. Our evaluation concerns transaction-class
inference (separating read from select by CLRT) on one device, and the replacement of that device's
evaluated features; we do not claim device-identity results or indistinguishability among devices. The
adversary knows the protocol and may know the defense policy; its advantage comes only from observing
the features.

\subsection{Trusted computing base and deployment assumptions}
The master, the switch and its loaded program, and the outstation are trusted. The defense runs
entirely in the switch data plane, with no host agent, controller fast path, or endpoint change. Two
deployment assumptions bound the timing argument. First, the protected flow carries no TCP timestamp
option, which would be an independent timing channel; in our evaluation the option was disabled on the
host and confirmed absent at the packet level in a preflight check, and the switch program also parses
the option so the timing argument is asserted only when it is absent. Second, only the eligible
response class (here the 49-byte read and control responses) is shaped; other traffic is forwarded
unchanged.

\subsection{Security goals}
We state the security goals as testable properties of the master-facing observation. \textbf{G1:} the
defended CLRT is a single policy value independent of the transaction class. \textbf{G2:} every
eligible response presents the fixed segment shape $[28,21]$. \textbf{G3:} the master reassembles a
sequence-contiguous, CRC-valid, checksum-valid 49-byte response. \textbf{G4:} for a control command,
the master-visible acknowledgment and echo placement does not vary with the internal release delay.
\textbf{G5 (safety):} the defense actuates no physical output. A separate evaluation-validity
property, that native and defended traffic differ only by a runtime mode on one loaded program, lets
the comparison isolate the defense; we treat it as a methodology control rather than a security goal.

\subsection{Non-goals and evidence limitations}
The following are out of scope, and the evidence does not establish them. We do not verify the
relay-facing link: the internal release time and the number of copies delivered to the relay are not
captured, so the single release is a modeled, not a measured, property, and we do not claim
exactly-once delivery. We do not hide the total response length, which stays 49 bytes and is
recoverable by a full TCP reader. We do not claim byte identity between an emitted response and a
source frame, because no paired source-side oracle was captured. We do not generalize from one
SEL-751 and one eligible class to all DNP3 outstations, and we do not claim indistinguishability among
devices or resistance to every traffic-analysis feature. We do not defend against an active adversary
or add payload confidentiality. If the observer could also see the relay-facing link, it could time
the internal release directly, and the control-command timing property would no longer hold; the
defense is scoped to the master-facing observer.
